% Part V -- Geospatial Data and Applied Mapping
% Chapter 15: geodatasets and applied mapping

\h{GeoDatasets and Applied Mapping}
\label{ch:15}

Small hand-written coordinates are useful for learning, but applied spatial
analysis needs documented data and repeatable acquisition. The
\c{geodatasets} package provides a catalog and downloader for educational
spatial datasets. GeoPandas then reads, analyzes, joins, and maps those files.
This chapter combines Chicago community polygons with grocery locations to
demonstrate a complete spatial workflow.

\begin{NoteBox}
\textbf{Prerequisites.} Complete \booklink{ch:14}{Chapter 14} first. You should
understand GeoDataFrames, geometry, CRS transformation, pandas grouping, and
Matplotlib subplots. The first dataset request may need an internet connection;
later runs can use the cached local copy.
\end{NoteBox}

\begin{tcolorbox}[title=Learning outcomes,colframe=orange,colback=white,arc=2mm]
After completing this chapter, you should be able to:
\begin{itemize}
    \item distinguish a dataset catalog from a spatial-analysis library;
    \item retrieve a documented dataset path with \c{geodatasets.get_path()};
    \item load vector files with \c{gpd.read_file()} and inspect their schema;
    \item align layers in one projected CRS;
    \item build a choropleth and overlay a point layer;
    \item use a spatial join to attach polygon attributes to points;
    \item aggregate joined records and place a map beside a comparison chart;
          and
    \item state important limits of a descriptive spatial result.
\end{itemize}
\end{tcolorbox}

\hh{Know what each library contributes}

The packages have separate responsibilities. \c{geodatasets} knows dataset
names, source URLs, and local cache locations. GeoPandas reads spatial files
and performs spatial operations. Matplotlib controls the final Figure and Axes.

\begin{SyntaxBox}{Catalog, analysis, and presentation are separate stages}
\begin{itemize}
    \item \c{geodatasets.get_path(name)} obtains a usable local path, downloading
          and caching the source when necessary.
    \item \c{gpd.read_file(path)} converts the file into a GeoDataFrame.
    \item GeoPandas methods transform, join, and plot the spatial records.
    \item Matplotlib coordinates multiple panels, labels, and export.
\end{itemize}
Separating these roles makes failures easier to diagnose: acquisition, data
interpretation, analysis, and presentation are not the same problem.
\end{SyntaxBox}

\begin{ExplainBox}{A dataset is evidence with provenance, not just a file}
A defensible spatial result records who produced the data, what one feature
represents, when and where it was collected, its coordinate reference system,
important field definitions, and known omissions. A successful download proves
only that bytes were obtained. Inspection and metadata are needed before those
bytes can support a spatial claim. Cache paths improve reproducibility, but
they do not guarantee that a source is complete, current, or suitable for the
question.
\end{ExplainBox}

The package API and catalog are documented in the
\href{https://geodatasets.readthedocs.io/en/latest/api.html}{official geodatasets API reference}.

\hh{Discover, retrieve, and load a dataset}

Use a catalog name rather than embedding an unstable download URL in every
script. The following example requests Chicago community areas with population
attributes.

\begin{lstlisting}[
    language=python,
    style=block,
    caption={Retrieve and inspect an educational spatial dataset},
    label={c15_get_path},
    captionpos=b
]
import geopandas as gpd
import geodatasets

dataset_name = "geoda.chicago_commpop"

print(geodatasets.get_url(dataset_name))
dataset_path = geodatasets.get_path(dataset_name)
chicago = gpd.read_file(dataset_path)

print(chicago.shape)
print(chicago.columns.tolist())
print(chicago.crs)
print(chicago.geometry.geom_type.value_counts())
\end{lstlisting}

\begin{SyntaxBox}{Follow the data from identifier to GeoDataFrame}
\begin{enumerate}
    \item \c{dataset_name} is a catalog key, not a file path.
    \item \c{get_url()} reveals the external source for provenance checking.
    \item \c{get_path()} returns a local path and may download the file once.
    \item \c{gpd.read_file()} parses attributes, geometry, and CRS into one object.
    \item The final print statements verify dimensions, schema, CRS, and geometry.
\end{enumerate}
Cache behavior improves reproducibility, but it does not remove the obligation
to record the dataset name, source, fields used, and analysis date.
\end{SyntaxBox}

At the time this example was prepared, the layer contained community names in
\c{community}, 2010 population in \c{POP2010}, and Polygon or MultiPolygon
geometry. Always inspect the downloaded version instead of assuming that a
field name or record count can never change.

\hh{Load a second layer and align both CRS values}

The grocery dataset contains Point locations. Load both layers, then transform
them to the same projected CRS before overlaying or joining them.

\begin{lstlisting}[
    language=python,
    style=block,
    caption={Load community polygons and grocery points in one CRS},
    label={c15_align_layers},
    captionpos=b
]
chicago = gpd.read_file(
    geodatasets.get_path("geoda.chicago_commpop")
)
groceries = gpd.read_file(
    geodatasets.get_path("geoda.groceries")
)

analysis_crs = "EPSG:26916"
chicago = chicago.to_crs(analysis_crs)
groceries = groceries.to_crs(chicago.crs)

print(chicago.crs == groceries.crs)
print(len(chicago), len(groceries))
\end{lstlisting}

The second \c{to_crs()} reuses \c{chicago.crs}, so the two layers cannot drift
apart because of a repeated string typed differently. EPSG:26916 is a
meter-based UTM CRS suitable for this Chicago-area demonstration. A different
study location may require a different projected CRS.

\begin{TryBox}{Check alignment before plotting}
Print \c{chicago.total_bounds} and \c{groceries.total_bounds}. The grocery
extent should fall broadly inside the community extent, and both arrays should
have coordinates of a similar magnitude. Matching CRS labels alone do not
prove that the source data were declared correctly.
\end{TryBox}

\hh{Build a choropleth and point overlay}

A choropleth assigns polygon color from a numeric attribute. It is effective
for showing geographic variation, but its legend, variable definition, and
boundary context must be clear.

\begin{lstlisting}[
    language=python,
    style=block,
    caption={Map population and grocery locations together},
    label={c15_choropleth},
    captionpos=b
]
import matplotlib.pyplot as plt

fig, ax = plt.subplots(figsize=(7, 7), layout="constrained")

chicago.plot(
    column="POP2010",
    cmap="Blues",
    legend=True,
    edgecolor="white",
    linewidth=0.35,
    legend_kwds={"label": "Population in 2010", "shrink": 0.75},
    ax=ax,
)
groceries.plot(
    ax=ax,
    color="tab:orange",
    edgecolor="black",
    linewidth=0.25,
    markersize=10,
    zorder=3,
)

ax.set_title("Population and mapped grocery locations")
ax.set_axis_off()
plt.show()
plt.close(fig)
\end{lstlisting}

\begin{SyntaxBox}{Separate data mapping from map decoration}
\c{column="POP2010"} connects data values to polygon color, while
\c{cmap="Blues"} chooses one ordered color scale. \c{legend=True} exposes the
numeric scale. Polygon edges remain thin and neutral so the orange grocery
points retain emphasis. Both calls receive the same \c{ax}; \c{zorder=3} keeps
the smaller point symbols visible above the polygons.
\end{SyntaxBox}

The
\href{https://geopandas.org/en/stable/docs/user_guide/mapping.html}{official GeoPandas mapping guide}
provides additional parameters and classification options.

\hh{Attach community names with a spatial join}

The grocery table does not need to contain a community name already. A spatial
join can determine which community Polygon contains each grocery Point and
copy the selected polygon attributes to the point rows.

\begin{lstlisting}[
    language=python,
    style=block,
    caption={Join each grocery point to its containing community},
    label={c15_spatial_join},
    captionpos=b
]
joined = gpd.sjoin(
    groceries,
    chicago[["community", "geometry"]],
    how="left",
    predicate="within",
)

print(joined[["community", "geometry"]].head())
print("Unmatched points:", joined["community"].isna().sum())

counts = (
    joined.dropna(subset=["community"])
    .groupby("community")
    .size()
    .sort_values(ascending=False)
)
print(counts.head(10))
\end{lstlisting}

\begin{SyntaxBox}{Read the join and aggregation precisely}
\begin{itemize}
    \item The left object is the point layer whose rows should be retained.
    \item The right object supplies only \c{community} and \c{geometry}, reducing
          unnecessary duplicate columns.
    \item \c{how="left"} keeps every grocery row, including unmatched points.
    \item \c{predicate="within"} matches a Point when it lies inside a community
          geometry.
    \item \c{dropna()}, \c{groupby()}, and \c{size()} then count matched points by
          community.
\end{itemize}
Inspect unmatched rows before dropping them. A point can be unmatched because
of boundary precision, incomplete coverage, or an incorrect CRS.
\end{SyntaxBox}

The full method contract appears in the
\href{https://geopandas.org/en/stable/docs/reference/api/geopandas.sjoin.html}{official GeoPandas spatial-join reference}.

\hh{Applied example: combine a map and a ranked chart}
\label{sec:c15-application}

The final application integrates acquisition, inspection, CRS alignment,
spatial joining, aggregation, visualization, and export. One panel preserves
geographic context; the other makes the largest counts easier to compare.

A runnable copy is provided in \c{examples/c15_chicago_mapping.py}.

\begin{lstlisting}[
    language=python,
    style=block,
    caption={Create a Chicago population and grocery-location analysis},
    label={c15_chicago_application},
    captionpos=b
]
from pathlib import Path

import geopandas as gpd
import geodatasets
import matplotlib.pyplot as plt

# 1. Retrieve and load the two documented layers.
community_path = geodatasets.get_path("geoda.chicago_commpop")
grocery_path = geodatasets.get_path("geoda.groceries")

chicago = gpd.read_file(community_path)
groceries = gpd.read_file(grocery_path)

# 2. Inspect the minimum assumptions used later.
required_columns = {"community", "POP2010", "geometry"}
missing = required_columns.difference(chicago.columns)
if missing:
    raise ValueError(f"Missing Chicago columns: {sorted(missing)}")
if chicago.crs is None or groceries.crs is None:
    raise ValueError("Both layers must have a CRS")

# 3. Align the layers in a projected Chicago-area CRS.
chicago = chicago.to_crs("EPSG:26916")
groceries = groceries.to_crs(chicago.crs)

# 4. Attach a community name to each grocery point.
joined = gpd.sjoin(
    groceries,
    chicago[["community", "geometry"]],
    how="left",
    predicate="within",
)

# 5. Calculate the ten largest mapped-location counts.
top_counts = (
    joined.dropna(subset=["community"])
    .groupby("community")
    .size()
    .sort_values(ascending=False)
    .head(10)
    .sort_values()
)

# 6. Coordinate a map and a comparison chart.
fig, axes = plt.subplots(
    1,
    2,
    figsize=(10.8, 5.4),
    layout="constrained",
)

chicago.plot(
    column="POP2010",
    cmap="Blues",
    legend=True,
    edgecolor="white",
    linewidth=0.35,
    legend_kwds={"label": "Population in 2010", "shrink": 0.72},
    ax=axes[0],
)
groceries.plot(
    ax=axes[0],
    color="tab:orange",
    edgecolor="black",
    linewidth=0.25,
    markersize=9,
    zorder=3,
)
axes[0].set_title("Population and grocery locations")
axes[0].set_axis_off()

top_counts.plot.barh(ax=axes[1], color="tab:blue")
axes[1].set(
    title="Communities with the most mapped groceries",
    xlabel="Mapped grocery locations",
    ylabel="Community",
)
axes[1].grid(axis="x", alpha=0.2)

# 7. Record the source and export the complete Figure.
fig.text(
    0.5,
    0.01,
    "Source: GeoDa Center datasets distributed through geodatasets",
    ha="center",
    fontsize=8,
)
output_folder = Path("study_figures")
output_folder.mkdir(exist_ok=True)
fig.savefig(
    output_folder / "chicago_population_groceries.png",
    dpi=200,
    bbox_inches="tight",
)
plt.show()
plt.close(fig)
\end{lstlisting}

\bookfigure[0.98\textwidth]{Figures/Generated/c15_chicago_mapping.pdf}
{Rendered application: geographic context on the left and a ranked comparison of joined point counts on the right.}
{fig:c15-chicago-mapping}

\begin{SyntaxBox}{Trace the complete application by responsibility}
\begin{enumerate}
    \item \textbf{Acquire:} catalog names resolve to cached local paths.
    \item \textbf{Load and validate:} both GeoDataFrames and required fields are
          checked before analysis.
    \item \textbf{Transform:} a common projected CRS makes the geometries align.
    \item \textbf{Relate:} \c{sjoin()} transfers a community name to each point.
    \item \textbf{Summarize:} a pandas pipeline counts, ranks, and orders matches.
    \item \textbf{Communicate:} the map answers \emph{where}; the bar chart answers
          \emph{which communities have the largest mapped counts}.
    \item \textbf{Document:} a source note and deterministic file name accompany
          the exported Figure.
\end{enumerate}
Validation belongs inside the script because a visually plausible map can still
be based on missing columns, absent CRS metadata, or unmatched features.
\end{SyntaxBox}

\hh{Interpret responsibly}

The right panel counts mapped grocery records within community boundaries. It
does \emph{not} measure grocery access, store quality, food affordability, or
locations per resident. A larger community or a more populous community may
naturally contain more records. A defensible analysis should consider area or
population normalization, travel networks, dataset vintage, missing locations,
and the modifiable areal unit problem before making policy claims.

\begin{ExplainBox}{A map is an argument, not just an image}
Polygon boundaries define aggregation units; color encodes one selected
variable; symbol size and order direct attention; and missing records influence
the apparent pattern. State these choices and limitations so readers can
evaluate the evidence rather than treating the map as neutral.
\end{ExplainBox}

\hh{Common mistakes and useful diagnoses}

\begin{itemize}
    \item \textbf{The first run cannot obtain data.} Confirm internet access and
          print \c{geodatasets.get_url(name)} to identify the source.
    \item \textbf{A field name raises \c{KeyError}.} Print \c{columns.tolist()}
          and inspect the actual schema before analysis.
    \item \textbf{A spatial join returns many missing values.} Check both CRS
          values, extents, geometry validity, and predicate direction.
    \item \textbf{Points disappear below polygons.} Plot polygons first, points
          second, and assign the points a higher \c{zorder}.
    \item \textbf{Counts are interpreted as rates.} Choose a defensible
          denominator and explain what the normalized measure means.
    \item \textbf{A map has no provenance.} Record dataset names, source notes,
          relevant dates, variables, transformations, and exclusions.
\end{itemize}

\hh{Practice ladder}

\hhh{Level 0 -- Read}

Identify which four lines in Listing~\ref{c15_get_path} perform discovery,
retrieval, loading, and inspection.

\hhh{Level 1 -- Modify}

Change the choropleth from \c{POP2010} to \c{POPCH}. Update the legend and title
so a reader can understand what the field represents after inspecting its data.

\hhh{Level 2 -- Explain}

Explain why \c{how="left"} is preferable to \c{how="inner"} when diagnosing
unmatched grocery points.

\hhh{Level 3 -- Build}

Join grocery points to communities and calculate
\c{groceries_per_10000 = count / POP2010 * 10000}. Exclude or flag any community
with a nonpositive or missing population before division.

\hhh{Level 4 -- Challenge}

Write a function \c{spatial_count(points, polygons, group_column)} that checks
CRS compatibility, performs a left spatial join, reports unmatched points, and
returns a sorted count Series without mutating either input.

\hh{Practice solutions}

\textbf{Level 0:} \c{get_url()} discovers the source, \c{get_path()} retrieves a
local path, \c{read_file()} loads it, and the print calls inspect the result.

\textbf{Level 2:} A left join retains unmatched input points with missing polygon
attributes, making the problem visible. An inner join would silently remove them.

\textbf{Level 3}

\begin{lstlisting}[
    language=python,
    style=block,
    caption={One population-normalized comparison},
    label={c15_solution_level3},
    captionpos=b
]
counts = (
    joined.dropna(subset=["community"])
    .groupby("community")
    .size()
    .rename("grocery_count")
)

rates = chicago[["community", "POP2010"]].merge(
    counts,
    on="community",
    how="left",
)
rates["grocery_count"] = rates["grocery_count"].fillna(0)
rates = rates[rates["POP2010"] > 0].copy()
rates["groceries_per_10000"] = (
    rates["grocery_count"] / rates["POP2010"] * 10_000
)
print(rates.sort_values("groceries_per_10000", ascending=False).head())
\end{lstlisting}

This rate adjusts only for population count. It still does not represent travel
time, store capacity, price, or individual access.

\hh{Chapter checkpoint}

\begin{enumerate}
    \item What role does \c{geodatasets} play in the workflow?
    \item Which function converts a spatial file path into a GeoDataFrame?
    \item Why transform both layers to a shared CRS before a spatial join?
    \item What does \c{predicate="within"} mean when points are on the left?
    \item Why inspect unmatched rows after a left join?
    \item What question does a choropleth answer differently from a bar chart?
    \item Why is a mapped grocery count not automatically an access measure?
\end{enumerate}

\hh{Checkpoint answers}

\begin{enumerate}
    \item It catalogs, downloads, and caches documented educational datasets.
    \item \c{gpd.read_file()}.
    \item Spatial relations require geometry expressed in a consistent coordinate system.
    \item Each point is matched to a polygon that spatially contains it.
    \item Unmatched records may reveal CRS, coverage, boundary, or data-quality issues.
    \item A choropleth shows geographic variation; a bar chart supports direct ranking.
    \item Count ignores population, area, travel, price, quality, and missing records.
\end{enumerate}

\begin{tcolorbox}[title=Part V summary,colframe=orange,colback=white,arc=2mm]
GeoPandas connects table reasoning to geometry, CRS, spatial predicates, and
layered maps. \c{geodatasets} supports reproducible access to teaching data,
while spatial joins connect observations through location. A sound applied
workflow acquires, inspects, transforms, validates, analyzes, visualizes, and
qualifies its conclusions in that order.
\end{tcolorbox}
